\begin{textsample}{POS Dim 3 – ac – Score 41.00 – ac/ac\_inf\_21.txt}  \label{ex:f3_pos_026}
21. \textbf{CRIANÇAS} BEM \textbf{PEQUENAS} ( 1 ANO E 7 \textbf{MESES} A 3 ANOS E 11 \textbf{MESES} ) 21.1.
Campo de experiências “ Espaços, tempos, quantidades, relações e transformações ” As \textbf{crianças} vivem inseridas em espaços e tempos de diferentes dimensões, em um mundo constituído de fenômenos naturais e socioculturais.
Desde muito \textbf{pequenas}, elas procuram se situar em diversos espaços ( rua, bairro, cidade etc. ) e tempos ( dia e noite ; hoje, ontem e amanhã etc. ). \textbf{Demonstram} também \textbf{curiosidade} sobre o mundo físico ( seu próprio corpo, os fenômenos atmosféricos, os animais, as plantas, as transformações da natureza, os diferentes tipos de materiais e as possibilidades de sua manipulação etc. ) e o mundo sociocultural ( as relações de parentesco e sociais entre as pessoas que conhece ; como vivem e em que trabalham essas pessoas ; quais suas tradições e seus costumes ; a diversidade entre elas etc. ).
Além disso, nessas experiências e em muitas outras, as \textbf{crianças} também se deparam, \textbf{frequentemente}, com conhecimentos matemáticos ( contagem, ordenação, relações entre quantidades, dimensões, medidas, comparação de \textbf{pesos} e de comprimentos, avaliação de distâncias, reconhecimento de formas geométricas, conhecimento e reconhecimento de numerais cardinais e ordinais etc. ) que igualmente aguçam a \textbf{curiosidade}.
Portanto, a Educação \textbf{Infantil} precisa promover experiências nas quais as \textbf{crianças} possam fazer observações, \textbf{manipular} objetos, investigar e explorar seu entorno, levantar hipóteses e consultar fontes de informação para buscar respostas às suas \textbf{curiosidades} e indagações.
Assim, a instituição escolar está criando oportunidades para que as \textbf{crianças} ampliem seus conhecimentos do mundo físico e sociocultural e possam utilizá -los em seu cotidiano. 21.2.
Direitos de aprendizagem e desenvolvimento - Conviver com seus pares e com \textbf{adultos}, explorando objetos e materiais de diferentes propriedades físicas, fazendo observações e construindo hipóteses sobre o mundo natural e social. - \textbf{Brincar} com seus pares e \textbf{adultos}, utilizando objetos, acessórios e elementos da natureza, assumindo diferentes papéis sociais e culturais, ampliando sua experiência relacional e \textbf{sensorial} com \textbf{texturas}, cheiros, cores, tamanhos, \textbf{pesos}, densidades e possibilidades de transformação. - Participar de situações de \textbf{cuidados} com o meio ambiente e sua sustentabilidade, de investigação sobre fenômenos naturais e resolução de problemas cotidianos que envolvam quantidades, medidas, dimensões, tempos, espaços, comparações, transformações, construindo hipóteses e explorando objetos e ferramentas diversas, tais como : bússola, lanterna, lupa, máquina \textbf{fotográfica}, smartphone, filmadora, gravador, projetor, computador e outras. - Explorar características do mundo natural e social, despertando a \textbf{curiosidade} por nomear, agrupar e ordenar informações para compreender os espaços à sua \textbf{volta} e as transformações neles realizados, sua própria história, os modos de vida das pessoas de sua comunidade e de outras culturas. - Expressar por meio das diferentes linguagens, suas impressões, observações, hipóteses e explicações sobre acontecimentos sociais, fenômenos da natureza e características do ambiente em que vive. - Conhecer -se e construir sua identidade pessoal e cultural, identificando seus próprios interesses na relação com o mundo físico e social, convivendo e conhecendo os costumes, as crenças e as tradições locais, regionais e de outras culturas. 21.2.1.
Objetivo de aprendizagem e desenvolvimento : ( EI02ET01 ) - Explorar e descrever semelhanças e diferenças entre as características e propriedades dos objetos ( \textbf{textura}, massa, tamanho ).
Aprendizagens específicas - Desenvolver a \textbf{curiosidade} explorando o meio, fazendo descobertas quanto as características e propriedades dos objetos. - Desenvolver a postura investigativa, a partir de observações das características e propriedades dos objetos de diferentes formas, \textbf{textura}, massa, tamanho. - Desenvolver o interesse em explorar objetos identificando suas semelhanças e diferenças. - Manifestar seus conhecimentos, criando narrativas, nomeando as características dos objetos.
Sugestões de experiências - Explorar os objetos percebendo seus atributos ( massa, \textbf{textura} e tamanho ). - Descrever objetos em situações de exploração. - Explicar suas ações e interferências sobre o meio ( parar uma bola, fazer bolinhas de \textbf{areia}, etc. ). - \textbf{Manipular} elementos do meio, misturando e observando suas transformações ( água e maisena, tintas e outras ). - Participar de experiências culinárias observando as transformações dos ingredientes usados. - Explorar traços e formas utilizando materiais de \textbf{pintura}, escultura e dobradura. - \textbf{Brincar} com diferentes tipos de \textbf{tinta}. - Participar de experiências de confecção de diferentes tipos de tinturas. - \textbf{Brincar} com elementos da natureza – argila, \textbf{areia}, água e plantas. 21.2.2.
Objetivo de aprendizagem e desenvolvimento : ( EI02ET02 ) - Observar, relatar e descrever incidentes do cotidiano e fenômenos naturais ( luz solar, vento, chuva etc. ).
Aprendizagens específicas - Desenvolver a \textbf{curiosidade} e interesse em observar fenômenos da natureza. - Construir, gradativamente, saberes e conhecimentos sobre fenômenos da natureza. - Desenvolver estratégias investigativas, questionando, relacionando informações, formulando hipóteses sobre a complexidade e diversidade da natureza. - Ampliar seus conhecimentos sobre fenômenos da natureza conhecidos por meio da \textbf{escuta} de histórias e meios de comunicação. - Manifestar interesse em compartilhar, descrever, \textbf{demonstrar}, relatar oralmente seus conhecimentos e experiências sobre as transformações da natureza, dos objetos e dos seres vivos. - Desenvolver a oralidade, explicando \textbf{livremente} sua compreensão sobre os fenômenos da natureza. - Construir conhecimentos a partir de investigações para descobrir porque as coisas acontecem e como funcionam. - \textbf{Demonstrar} conhecimentos sobre os movimentos do sol, da lua, das estrelas, das nuvens, e de outros elementos da natureza, a partir da observação investigativa. - Expressar oralmente suas impressões sobre o clima : calor, chuva, claro-escuro, quente-frio.
Sugestões de experiências - Investigar para descobrir porque as coisas acontecem e como funcionam. - Falar sobre o fenômeno natural que presenciou e/ou que está presenciando. - Descrever mudanças em objetos, seres vivos e eventos naturais no ambiente interno e na natureza. - Observar os diferentes elementos e fenômenos da natureza – luz solar, chuva, vento, nuvens, morros, dentre outros. - Observar e \textbf{conversar} sobre o movimento da lua, das estrelas e das nuvens. - \textbf{Contar} suas observações sobre as mudanças de tempo – frio e calor – nas diversas situações cotidianas. 21.2.3.
Objetivo de aprendizagem e desenvolvimento : ( EI02ET03 ) - Compartilhar, com outras \textbf{crianças}, situações de \textbf{cuidado} de plantas e animais nos espaços da instituição e fora dela.
Aprendizagens específicas - Desenvolver atitudes de \textbf{cuidados} com o ambiente da instituição educativa e outros espaços de convívio social. - Desenvolver atitude de preservação da natureza, percebendo -se como parte integrante e atuante do meio ambiente. - Desenvolver atitude de respeito, \textbf{cuidado} e permanente interesse por aprender sobre plantas e animais e suas relações com o ambiente. - Ampliar suas noções e compreensões sobre plantas e animais, refletindo sobre suas relações com o ambiente. - Manifestar conhecimentos sobre algumas particularidades e características dos animais, nomeando, \textbf{imitando}, cuidando, descrevendo e registrando por meio de diferentes linguagens. - Manifestar conhecimentos sobre as plantas e suas características, nomeando, cuidando, descrevendo e registrando por meio de diferentes linguagens. - Desenvolver o prazer da descoberta, por meio da exploração, observação, formulação de hipóteses sobre a diferença entre seres vivos e outros elementos e materiais.
Sugestões de experiências - Cuidar das plantas utilizando ferramentas para jardinagem ou horta. - Acompanhar o crescimento de alimentos na horta. - \textbf{Imitar} algumas particularidades dos animais. - Citar características que diferenciam os seres vivos de outros elementos. - Observar animais em livros, revistas e filmes. - Reproduzir \textbf{sons} de animais. - Descrever características físicas de animais : pelagem, forma do corpo, presença de bico, localização dos olhos e outras. - Relatar sobre diferentes formas de vida dos animais : alimentação, moradia, dentre outras. 21.2.4.
Objetivo de aprendizagem e desenvolvimento : ( EI02ET04 ) - Identificar relações espaciais ( dentro e fora, em cima, embaixo, acima, abaixo, entre e do lado ) e temporais ( antes, durante e depois ).
Aprendizagens específicas - Desenvolver noção de tempo e de espaço por meio de estratégias de observação e uso no cotidiano. - Perceber a passagem do tempo, acompanhando a retomada do professor das diferentes atividades da \textbf{rotina} \textbf{diária}. - Perceber algumas passagens significativas do tempo, participando da organização de eventos, festas tradicionais e aniversários. - Manifestar, em conversa com seus pares e \textbf{adultos}, conhecimentos sobre as relações temporais usando as expressões antes, durante e depois. - Desenvolver noções de espaço e tempo, tendo primeiramente seu corpo e suas ações como referências em \textbf{brincadeiras} livres. - Desenvolver atitudes de interesse por conhecer diferentes espaços da instituição educativa, por meio de explorações e deslocamentos. - Manifestar conhecimentos sobre percursos e trajetos considerando diferentes pontos de referência.
Sugestões de experiências - Procurar objetos ou \textbf{brinquedos} desejados nas situações de \textbf{brincadeiras}. - Usar expressões temporais como antes, durante e depois, em diferentes momentos da \textbf{rotina}. - Explorar os diferentes espaços da escola. - Descrever, percursos e trajetos considerando diferentes pontos de referência. - Consultar o calendário em diferentes situações. - \textbf{Brincar} de localizar um objeto no espaço. - \textbf{Brincar} de circuitos com diferentes obstáculos. - \textbf{Brincar} de faz-de-conta organizando o espaço. 21.2.5.
Objetivo de aprendizagem e desenvolvimento : ( EI02ET05 ) - Classificar objetos, considerando determinado atributo ( tamanho, \textbf{peso}, cor, forma etc. ).
Aprendizagens específicas - Desenvolver, \textbf{progressivamente}, a habilidade de classificar os objetos baseados em seus atributos : tamanho, \textbf{peso}, cor, forma e outros atributos. - Aprimorar os conhecimentos sobre os atributos dos objetos, em \textbf{brincadeiras} e explorações. - Manifestar suas descobertas sobre os atributos dos objetos, identificando, ordenando, selecionando, agrupando, nomeando e organizando as informações por meio de suas ações. - Desenvolver a capacidade investigativa, explorando, comparando as diferenças entre os materiais e classificando -os quanto ao tamanho, \textbf{peso}, cor, forma, dentre outros. - Desenvolver \textbf{hábito} de organizar os \textbf{brinquedos} e outros materiais da sala, selecionando e agrupando -os segundo seus atributos.
Sugestões de experiências - \textbf{Brincar} com materiais de tamanho, \textbf{peso}, cor e formato diferentes. - Organizar os \textbf{brinquedos} da sala, agrupando -os conforme seus atributos. - \textbf{Brincar} com diferentes instrumentos de medidas. - \textbf{Brincar} com o corpo observando os diferentes tamanhos ( braços, pernas, mãos, \textbf{pés}, etc. ). - \textbf{Brincar} de \textbf{montar} cenários com objetos de diversos tamanhos. - \textbf{Brincar} dentro de caixas e de tendas de tamanhos diversos. - Pintar usando suportes de diversos tamanhos e formatos. - Explorar \textbf{brinquedos} com diferentes formas. - \textbf{Brincar} com blocos de \textbf{encaixe}. 21.2.6.
Objetivo de aprendizagem e desenvolvimento : ( EI02ET06 ) - Utilizar conceitos básicos de tempo ( agora, antes, durante, depois, ontem, hoje, amanhã, lento, rápido, depressa, devagar ).
Aprendizagens específicas - Desenvolver, \textbf{progressivamente}, a capacidade de compreensão dos conceitos básicos de tempo ( agora, antes, durante, depois, ontem, hoje, amanhã, lento, rápido, depressa, devagar ) nas diferentes situações do cotidiano. - Desenvolver a percepção do tempo, tendo como referência os diferentes momentos da \textbf{rotina}, calendário, relógio e outras formas de perceber a passagem do tempo. - Reconhecer alguns indícios externos ( relacionados ao ambiente ) para antecipar acontecimentos do cotidiano. - Reconhecer, gradativamente, os diferentes níveis de velocidades e ritmos, as relações de causa e efeito dos diversos movimentos corporais, nas \textbf{brincadeiras} e outras experimentações. - Perceber algumas passagens significativas do tempo, participando da organização de eventos, festas tradicionais e aniversários.
Sugestões de experiências - \textbf{Brincar} nos diferentes espaços internos e externos da instituição. - \textbf{Brincar} experimentando diversos níveis de velocidade. - \textbf{Brincar} com jogos diversos. - \textbf{Conversar} sobre os diferentes momentos da \textbf{rotina}. - \textbf{Brincar} com relógio, calendário, fita métrica, balança e outros materiais convencionais e/ou não convencionais que envolvem medições. - \textbf{Brincar} com cantigas, rimas, parlendas, quadrinhas que se utilizam de números, contagens e medições. 21.2.7.
Objetivo de aprendizagem e desenvolvimento : ( EI02ET07 ) - \textbf{Contar} oralmente objetos, pessoas, livros etc., em contextos diversos.
Aprendizagens específicas - Construir, gradativamente, conceito de número em situações contextualizadas e significativas. - Compreender, \textbf{progressivamente}, o sistema numérico por meio de suas interações com as pessoas e com os materiais. - Apropriar -se, gradativamente, de estratégias de contagem da sequência numérica, nas \textbf{brincadeiras}, cantigas de roda, recitações numéricas, exploração de objetos e em outras situações. - Familiarizar -se com diferentes instrumentos ( calculadora, régua, fita métrica, teclado de computador, calendário etc. ) que possibilite usar e pensar sobre os números em contextos significativos.
Sugestões de experiências - Explorar os números naturais nas diferentes situações de uso cotidiano. - Recitar sequência numérica em \textbf{brincadeiras}. - \textbf{Brincar} explorando quantidades. - Participar de jogos de percurso. - \textbf{Brincar} de \textbf{contar} objetos em diferentes situações ( coleções, peças de jogos, cesto de tesouros, \textbf{brinquedos} etc ). - \textbf{Brincar} com objetos comparando quantidades. - \textbf{Brincar} de faz-de-conta com telefone, máquina de calcular, relógio e outros instrumentos. - \textbf{Contar} os \textbf{colegas} em situações de \textbf{brincadeiras}. - \textbf{Brincar} de juntar objetos ou \textbf{brinquedos} e distribuir para os \textbf{colegas}. - \textbf{Brincar} de repartir igualmente objetos. - \textbf{Brincar} de completar objetos de uma coleção. 21.2.8.
Objetivo de aprendizagem e desenvolvimento : ( EI02ET08 ) - Registrar com números a quantidade de \textbf{crianças} ( \textbf{meninas} e \textbf{meninos}, presentes e ausentes ) e a quantidade de objetos da mesma natureza ( bonecas, bolas, livros etc. ).
Aprendizagens específicas - Construir, mesmo que não-convencionalmente, formas de registrar números. - Representar quantidades por meio de desenhos, jogos, \textbf{brincadeiras} e manipulação de diferentes objetos. - Ampliar, \textbf{progressivamente}, a capacidade de realizar a leitura de números, ainda que não-convencionalmente, em \textbf{brincadeiras}, jogos cantados, parlendas, dentre outras possibilidades. - \textbf{Demonstrar} interesse e prazer por descobrir os números escritos em diferentes suportes. - Construir \textbf{pequenas} coleções que lhes sejam atraentes.
Sugestões de experiências - \textbf{Manusear} diferentes suportes que contenham números escritos. - Ler números escritos. - \textbf{Contar} e registrar quantidades, de forma não convencional, em jogos, \textbf{brincadeiras}, coleções e outras propostas. - Participar de produção de cartazes em que se registra números. - Trocar ideias com os \textbf{colegas} sobre os números. 21.3.
Avaliação - \textbf{acompanhamento} do processo de aprendizagem e desenvolvimento das \textbf{crianças} bem \textbf{pequenas} : - Observação \textbf{criteriosa}, \textbf{registro} e análise quanto : - À manifestação de conhecimentos construídos sobre as semelhanças e diferenças entre as características e propriedades dos objetos. - Ao interesse em compartilhar saberes e conhecimentos sobre fenômenos da natureza e incidentes do cotidiano. - À atitude de \textbf{cuidado} e interesse por aprender sobre plantas e animais e suas relações com o ambiente. - À manifestação de conhecimentos sobre as relações espaciais e temporais, em diferentes situações do cotidiano. - Às estratégias que elaboram para classificar objetos considerando seus atributos. - À compreensão dos conceitos básicos de tempo. - Às estratégias que utiliza para realizar contagens orais. - Às estratégias que utiliza para registrar números para representar quantidades.

% matched lemmas: acompanhamento, adulto, areia, brincadeira, brincar, brinquedo, colega, contar, conversar, criança, criterioso, cuidado, curiosidade, demonstrar, diário, encaixe, escuta, fotográfico, frequentemente, hábito, imitar, infantil, livremente, manipular, manusear, menino, montar, mês, pequeno, peso, pintura, progressivamente, pé, registro, rotina, sensorial, som, textura, tinta, volta
\end{textsample}
